%% 42장 최적화의 바탕
\chapter{Foundations of Optimization}
\label{ch:optimization-basics}

\abstract{통계의 최대우도 추정, 회귀의 최소제곱, 기계학습의 손실 최소화는 모두 함수의 최솟값을 찾는 문제이다. 이 권에서는 그 계산을 다룬다. 이 장은 뒤에 나올 알고리즘이 모두 기대는 이론을 정리한다. 먼저 극소점의 일계와 이계 조건을 증명한 뒤 볼록 집합과 볼록 함수의 정의와 특징을 정리하고 볼록 함수에서는 국소 최소가 그대로 전역 최소임을 보인다. 기울기가 Lipschitz 연속일 때 성립하는 하강 보조정리로 경사 하강법의 수렴 속도를 증명한다. 강볼록성이 더해지면 수렴이 선형이 되고 그 비율은 조건수가 정하는데 이 결과를 27장의 켤레 기울기법과 견준다. 끝으로 볼록 집합으로의 사영 정리를 증명하고 binary64에서 함수값만으로는 최솟값의 위치를 상대 정확도 $\sqrt u$보다 정확히 찾을 수 없는 까닭을 살핀다.}

\section{Optimality Conditions}
\label{sec:optimality}

\subsection{Local and Global Minima}

함수 $f:\mathbb R^n\to\mathbb R$에 대해 $x_\ast$의 어떤 근방에서 $f(x_\ast)\le f(x)$이면 $x_\ast$를 \term{국소 최소점}{local minimizer}, 모든 $x$에서 그러면 \term{전역 최소점}{global minimizer}이라 한다. 전역 최소점이 늘 있는 것은 아니다. $f$가 연속이고 하나의 준위 집합 $\{x:f(x)\le f(x_0)\}$이 유계이면, 그 준위 집합은 컴팩트하므로 Weierstrass 정리에 의해 최솟값이 있다. $\|x\|\to\infty$일 때 $f(x)\to\infty$인 함수를 \term{강제적}{coercive} 함수라 하며 강제적 함수는 이 조건을 만족한다.

\subsection{First- and Second-Order Conditions}

\begin{theorem}
\label{thm:optimality}
$f$가 $x_\ast$ 근처에서 두 번 연속 미분 가능하다고 하자.
\begin{enumerate}
\item $x_\ast$가 국소 최소점이면 $\nabla f(x_\ast)=0$이고 $\nabla^2f(x_\ast)$는 양반정치이다.
\item $\nabla f(x_\ast)=0$이고 $\nabla^2f(x_\ast)$가 양정치이면 $x_\ast$는 엄격한 국소 최소점이다.
\end{enumerate}
\end{theorem}

\begin{proof}
(i) 임의의 방향 $d$에 대해 $\phi(t)=f(x_\ast+td)$는 $t=0$에서 국소 최소이므로 $\phi'(0)=\nabla f(x_\ast)^{\mathsf T}d=0$이다. $d$가 임의이므로 $\nabla f(x_\ast)=0$이다. Taylor 정리로 $\phi(t)=\phi(0)+(t^2/2)d^{\mathsf T}\nabla^2f(x_\ast+\theta td)d$이고 작은 $t$에서 $\phi(t)\ge\phi(0)$이므로 $d^{\mathsf T}\nabla^2f(x_\ast+\theta td)d\ge0$이며, $t\to0$에서 연속성으로 $d^{\mathsf T}\nabla^2f(x_\ast)d\ge0$이다.

(ii) $\nabla^2f(x_\ast)$의 가장 작은 고유값을 $\lambda>0$이라 하면 연속성에 의해 $x_\ast$의 어떤 근방에서 $\nabla^2f$의 가장 작은 고유값이 $\lambda/2$ 이상이다. 그 근방의 $x\ne x_\ast$에서 Taylor 정리로 $f(x)=f(x_\ast)+2^{-1}(x-x_\ast)^{\mathsf T}\nabla^2f(\xi)(x-x_\ast)\ge f(x_\ast)+(\lambda/4)\|x-x_\ast\|^2>f(x_\ast)$이다.
\end{proof}

$\nabla f(x)=0$인 점을 \term{정류점}{stationary point}이라 한다. 정류점은 최소점일 수도, 최대점일 수도, 안장점일 수도 있다. 이 권의 알고리즘은 대개 정류점을 찾으며 찾은 점이 최소점인지는 함수의 구조나 이계 조건에 기대어 판단한다.

\section{Convexity}
\label{sec:convexity}

\subsection{Convex Sets and Convex Functions}

집합 $C$가 그 안의 두 점을 잇는 선분을 모두 포함하면, 곧 $x,y\in C$와 $t\in[0,1]$에서 $tx+(1-t)y\in C$이면 \term{볼록 집합}{convex set}이라 한다. 함수 $f$가 볼록 집합 위에서
\[
f\bigl(tx+(1-t)y\bigr)\le tf(x)+(1-t)f(y)
\]
를 만족하면 \term{볼록 함수}{convex function}라 한다. 그래프 위의 두 점을 이은 선분이 그래프보다 아래로 내려가지 않는다는 말이다. 최적화에서 볼록 함수를 따로 떼어 다루는 까닭은 다음 정리에서 드러난다\cite{boyd2004convex,rockafellar1970}.

\begin{theorem}
\label{thm:convex-characterization}
$f$가 열린 볼록 집합 위에서 미분 가능하면 다음은 동치이다. (i) $f$는 볼록이다. (ii) 모든 $x,y$에서 $f(y)\ge f(x)+\nabla f(x)^{\mathsf T}(y-x)$이다. $f$가 두 번 미분 가능하면 이것은 다시 (iii) 모든 $x$에서 $\nabla^2f(x)$가 양반정치인 것과 동치이다.
\end{theorem}

\begin{proof}
(i)$\Rightarrow$(ii): 볼록성에서 $f(x+t(y-x))\le f(x)+t(f(y)-f(x))$이므로 $\bigl(f(x+t(y-x))-f(x)\bigr)/t\le f(y)-f(x)$이고, $t\to0^+$으로 보내면 좌변이 $\nabla f(x)^{\mathsf T}(y-x)$이다. (ii)$\Rightarrow$(i): $z=tx+(1-t)y$에서 (ii)를 $x$와 $y$에 각각 쓰면 $f(x)\ge f(z)+\nabla f(z)^{\mathsf T}(x-z)$, $f(y)\ge f(z)+\nabla f(z)^{\mathsf T}(y-z)$이고, 두 식에 $t$와 $1-t$를 곱해 더하면 $t(x-z)+(1-t)(y-z)=0$이므로 $tf(x)+(1-t)f(y)\ge f(z)$이다. (i)$\Leftrightarrow$(iii)은 한 변수 함수 $\phi(s)=f(x+s(y-x))$로 줄여, 다음 사실을 쓰면 된다. 미분 가능한 한 변수 함수는 도함수가 단조 증가할 때 그리고 그때에만 볼록이고 도함수의 단조 증가는 이계 도함수가 음이 아닌 것과 같다.
\end{proof}

(ii)에 따르면 볼록 함수의 그래프는 어느 접평면보다도 아래로 내려가지 않는다. 그래서 정류점 $\nabla f(x_\ast)=0$에서 $f(y)\ge f(x_\ast)$가 모든 $y$에 대해 성립한다. 볼록 함수의 정류점은 전역 최소점이고 국소 최소점도 모두 전역 최소점이다. 최소점의 집합은 볼록 집합이다. 이 성질 덕분에 볼록 최적화에서는 찾은 답이 전역 최소인지 걱정할 필요가 없다. 22장의 최소제곱, 63장의 로지스틱 회귀의 음의 로그 우도, 64장의 lasso가 모두 볼록 문제이다.

\subsection{Strong Convexity and Smoothness}

수렴 속도를 말하려면 두 상수가 필요하다. $f(y)\ge f(x)+\nabla f(x)^{\mathsf T}(y-x)+(\mu/2)\|y-x\|^2$이 성립하는 $\mu>0$이 있으면 $f$를 $\mu$-\term{강볼록}{strongly convex}이라 하고, $\|\nabla f(x)-\nabla f(y)\|\le L\|x-y\|$이면 $f$를 $L$-\term{매끄럽다}{$L$-smooth}고 한다. $f$가 두 번 미분 가능하면 강볼록성은 $\nabla^2f\succeq\mu I$와 같고 $L$-매끄러움은 $-LI\preceq\nabla^2f\preceq LI$와 같다. $f$가 볼록이면 뒤의 조건은 $\nabla^2f\preceq LI$로 줄어든다. 비 $\kappa=L/\mu$를 문제의 \term{조건수}{condition number}라 한다. 이차 함수 $x^{\mathsf T}Ax/2-b^{\mathsf T}x$에서 $\mu$와 $L$은 $A$의 가장 작은 고유값과 가장 큰 고유값이고 $\kappa$는 4장의 조건수 $\kappa_2(A)$이다.

\begin{proposition}[하강 보조정리]
\label{prop:descent-lemma}
$f$가 $L$-매끄러우면 모든 $x,y$에서
\[
f(y)\le f(x)+\nabla f(x)^{\mathsf T}(y-x)+\frac L2\|y-x\|^2
\]
이다.
\end{proposition}

\begin{proof}
$f(y)-f(x)=\int_0^1\nabla f(x+t(y-x))^{\mathsf T}(y-x)\,dt$이므로 $f(y)-f(x)-\nabla f(x)^{\mathsf T}(y-x)=\int_0^1\bigl(\nabla f(x+t(y-x))-\nabla f(x)\bigr)^{\mathsf T}(y-x)\,dt$이다. Cauchy--Schwarz와 Lipschitz 조건으로 피적분함수는 $Lt\|y-x\|^2$ 이하이고, $\int_0^1t\,dt=1/2$이다.
\end{proof}

\section{Gradient Descent}
\label{sec:gradient-descent}

가장 단순한 알고리즘은 기울기의 반대 방향으로 가는 \term{경사 하강법}{gradient descent} $x_{k+1}=x_k-\alpha\nabla f(x_k)$이다. Cauchy가 1847년에 연립방정식을 풀려고 제안했다\cite{cauchy1847,lemarechal2012cauchy}. \cref{prop:descent-lemma}에 $y=x-\alpha\nabla f(x)$를 넣으면 $f(y)\le f(x)-\alpha(1-L\alpha/2)\|\nabla f(x)\|^2$이므로, $\alpha=1/L$에서 걸음마다
\begin{equation}
\label{eq:sufficient-decrease}
f(x_{k+1})\le f(x_k)-\frac1{2L}\|\nabla f(x_k)\|^2
\end{equation}
만큼 함수값이 준다.

\begin{theorem}
\label{thm:gd-rates}
$f$가 $L$-매끄러운 볼록 함수이고 최소점 $x_\ast$가 있으면, $\alpha=1/L$의 경사 하강법은
\[
f(x_k)-f(x_\ast)\le\frac{L\|x_0-x_\ast\|^2}{2k}
\]
를 만족한다. $f$가 $\mu$-강볼록이기도 하면 $f(x_k)-f(x_\ast)\le(1-\mu/L)^k\bigl(f(x_0)-f(x_\ast)\bigr)$이다.
\end{theorem}

\begin{proof}
볼록성에서 $f(x_k)-f(x_\ast)\le\nabla f(x_k)^{\mathsf T}(x_k-x_\ast)$이다. 이것과 \eqref{eq:sufficient-decrease}를 합치고 $g=\nabla f(x_k)$로 쓰면
\begin{align*}
f(x_{k+1})-f(x_\ast)&\le g^{\mathsf T}(x_k-x_\ast)-\frac1{2L}\|g\|^2=\frac L2\Bigl(\|x_k-x_\ast\|^2-\|x_k-g/L-x_\ast\|^2\Bigr)\\
&=\frac L2\Bigl(\|x_k-x_\ast\|^2-\|x_{k+1}-x_\ast\|^2\Bigr)
\end{align*}
이다. $k=0,\dots,K-1$에 대해 더하면 오른쪽이 망원 급수가 되어 $\sum_{k=1}^K(f(x_k)-f(x_\ast))\le(L/2)\|x_0-x_\ast\|^2$이다. \eqref{eq:sufficient-decrease}에 의해 $f(x_k)$가 감소하므로 좌변은 $K(f(x_K)-f(x_\ast))$ 이상이다.

강볼록이면 강볼록성의 정의에서 $y$에 대해 오른쪽을 최소화하면 $f(x_\ast)\ge f(x)-\|\nabla f(x)\|^2/(2\mu)$, 곧 $\|\nabla f(x)\|^2\ge2\mu\bigl(f(x)-f(x_\ast)\bigr)$이다. 이것을 \eqref{eq:sufficient-decrease}에 넣으면 $f(x_{k+1})-f(x_\ast)\le(1-\mu/L)\bigl(f(x_k)-f(x_\ast)\bigr)$이다.
\end{proof}

강볼록인 경우 비율이 $1-1/\kappa$이므로 오차를 $\epsilon$배로 줄이려면 대략 $\kappa\ln(1/\epsilon)$걸음이 든다. 27장의 켤레 기울기법은 이차 함수에서 $\sqrt\kappa$에 비례하는 걸음이면 충분했다. 고유값이 $[1,\kappa]$에 고르게 퍼진 $100\times100$ 이차 함수에서 함수값의 오차가 $10^{-10}$배로 줄 때까지 걸음 수를 세어 보았다. $\kappa=10$에서 경사 하강법이 $99$걸음, CG가 $18$걸음이고, $\kappa=100$에서 $1033$걸음 대 $41$걸음, $\kappa=1000$에서 $10963$걸음 대 $51$걸음이다. 경사 하강법의 걸음은 \cref{thm:gd-rates}의 상계 $219$, $2292$, $23015$의 대략 절반인데, 이차 함수에서는 함수값의 오차가 걸음마다 $(1-1/\kappa)^2$배로 줄기 때문이다. 경사 하강법은 걸음 하나가 싸지만 조건수가 나쁜 문제에서는 쓸모가 없다. 44장에서는 곡률 정보를 써서 이 의존성을 줄이는 방법을 다룬다\cite{nesterov2018}.

\section{Projection onto Convex Sets}
\label{sec:projection}

제약이 있는 최적화에서는 집합 밖으로 나간 점을 실현 가능 집합으로 되돌려야 한다.

\begin{theorem}[사영 정리]
\label{thm:projection}
$C\subset\mathbb R^n$이 공집합이 아닌 닫힌 볼록 집합이면 모든 $x$에 대해 $\|x-p\|$를 최소로 하는 $p\in C$가 하나로 존재한다. 이 점을 $P_C(x)$로 쓰면 $p=P_C(x)$일 필요충분조건은 모든 $y\in C$에서 $(x-p)^{\mathsf T}(y-p)\le0$인 것이다.
\end{theorem}

\begin{proof}
존재: $c_0\in C$를 하나 잡으면 $\{y\in C:\|x-y\|\le\|x-c_0\|\}$은 공집합이 아닌 컴팩트 집합이고 그 위에서 연속함수 $\|x-y\|$가 최솟값을 가진다. 특징: $p$가 최소점이고 $y\in C$이면 $t\in[0,1]$에서 $p+t(y-p)\in C$이므로 $\|x-p-t(y-p)\|^2\ge\|x-p\|^2$, 곧 $-2t(x-p)^{\mathsf T}(y-p)+t^2\|y-p\|^2\ge0$이고 $t\to0^+$에서 부등식을 얻는다. 거꾸로 부등식이 성립하면 $\|x-y\|^2=\|x-p\|^2-2(x-p)^{\mathsf T}(y-p)+\|y-p\|^2\ge\|x-p\|^2$이다. 유일성: $p$와 $q$가 모두 최소점이면 두 부등식 $(x-p)^{\mathsf T}(q-p)\le0$과 $(x-q)^{\mathsf T}(p-q)\le0$을 더해 $\|p-q\|^2\le0$이다.
\end{proof}

상자 $\{l\le x\le u\}$에 사영하려면 성분마다 잘라 내면 된다. 48장의 상자 제약 문제와 44장의 L-BFGS-B가 이 사영을 쓴다.

\section{Limits of Finite Precision}
\label{sec:optimization-precision}

최솟값 근처에서 함수는 평평하다. $f(x_\ast+h)\approx f(x_\ast)+f''(x_\ast)h^2/2$이므로, $|f''|h^2/2$이 $u|f(x_\ast)|$보다 작으면 binary64로 계산한 함수값은 $f(x_\ast)$와 구별되지 않는다. 그래서 함수값만 보고 최소점의 위치를 찾으면 그 정확도는 대략 $\sqrt{u|f|/|f''|}$에 머물고 상대 정확도로 치면 $\sqrt u\approx10^{-8}$ 수준이다. 예컨대 $f(x)=(x-1)^2+1$을 binary64로 계산하면 적어도 $|h|\le2^{-27}\approx7.5\times10^{-9}$인 모든 $x=1+h$에서 $f(x)$가 정확히 $1$이고, 실제로 그런 $h$는 $|h|\approx1.05\times10^{-8}$까지 이어진다. 33장에서 앞쪽 차분의 정확도가 $\sqrt u$에 머문 것과 모양이 닮았다. 거기서는 절단 오차와 반올림 오차의 균형이, 여기서는 최소점 근처의 이차적인 평평함이 제곱근을 낳는다. 기울기를 쓰면 사정이 낫다. $\nabla f(x_\ast+h)\approx\nabla^2f\,h$는 $h$에 비례하므로 기울기가 정확히 계산되는 한 위치를 $u$에 가깝게 정할 수 있다. 이 차이는 다음 장의 한 변수 최소화에서 자세히 다룬다.

\begin{instardust}[\code{docs/design/optimize.md} §1.1, §4]
stardust의 \code{sd.optimize}는 scipy의 개별 함수 이름(\code{fmin}, \code{fmin\_bfgs}, \code{brentq} 등)을 두지 않고 \code{minimize}, \code{minimize\_scalar}, \code{root}, \code{root\_scalar}, \code{least\_squares} 같은 통합 함수만 두어 방법을 \code{method}로 고르게 했다. 결과의 \code{converged}는 그 방법의 멈춤 판정을 만족했는지만 알려 줄 뿐 수학적인 최적성을 보장하지 않는다. 설계 문서는 계약을 판정의 뜻 그대로 적는다. 예컨대 BFGS는 $\|\nabla f\|_\infty\le\mathrm{gtol}$이다. 반복은 C++ 커널이 돌리고 사용자 함수는 Python 콜백으로 부른다. 콜백을 부르는 동안에는 GIL을 쥐고 있다가 무거운 선형대수 계산에서만 놓는다.
\end{instardust}

\section{Concluding Remarks}

최적화의 이론에는 두 줄기가 있다. 하나는 최소점에서 기울기와 Hessian이 만족하는 조건이고 다른 하나는 알고리즘이 그 점에 다가가는 속도인데, 이 속도는 매끄러움과 강볼록성이 정한다. 볼록성이 있으면 두 줄기가 모두 단순해진다. 국소 최소와 전역 최소의 구별이 사라지고, 강볼록성까지 있으면 수렴 속도는 조건수 하나로 요약된다. 다음 장은 가장 단순한 경우인 한 변수 함수의 최소화에서 출발해 도함수 없이 구간을 줄여 가는 방법과 그 정확도의 한계를 살핀다.

최적화의 표준 교과서는 Nocedal과 Wright\cite{nocedal2006}이고, 볼록 최적화는 Boyd와 Vandenberghe\cite{boyd2004convex}, 일차 방법의 이론은 Beck\cite{beck2017}이 자세하다.

\chapterreferences
